% transformer_diagram.tex
% Diagram transformerového modelu — dve fázy, konceptuálny
% Použitie: \resizebox{0.88\textwidth}{!}{% transformer_diagram.tex
% Diagram transformerového modelu — dve fázy, konceptuálny
% Použitie: \resizebox{0.88\textwidth}{!}{% transformer_diagram.tex
% Diagram transformerového modelu — dve fázy, konceptuálny
% Použitie: \resizebox{0.88\textwidth}{!}{% transformer_diagram.tex
% Diagram transformerového modelu — dve fázy, konceptuálny
% Použitie: \resizebox{0.88\textwidth}{!}{\input{figures/tikz/transformer_diagram}}
\begin{tikzpicture}[
    font=\small,
    box/.style={
        draw, rounded corners=5pt, align=center,
        minimum width=3.6cm, minimum height=0.85cm,
        inner sep=5pt
    },
    inputbox/.style={box, fill=gray!10, draw=gray!50},
    attnbox/.style={box, fill=orange!10, draw=orange!50},
    mergebox/.style={box, fill=purple!8, draw=purple!40},
    outbox/.style={box, fill=green!8, draw=green!50},
    labelbox/.style={font=\scriptsize\itshape, text=gray!65, align=left,
                     text width=5.0cm},
    arrow/.style={-{Latex[length=2.5mm]}, thick},
    dasharrow/.style={-{Latex[length=2mm]}, dashed, gray!60}
]

%% ══════════════════════════════════════════════════════════════════
%% FÁZA 1 — Kódovanie vrcholov
%% ══════════════════════════════════════════════════════════════════
\node[font=\bfseries\small, text=orange!70!black]
    at (0, 0.55) {Fáza 1 — kódovanie vrcholov};

\node[inputbox] (coords) at (0, -0.45)
    {Súradnice miest \((u_i, v_i)\)\\[-2pt]
     \scriptsize pre všetky vrcholy inštancie};

\node[attnbox] (embed) at (0, -2.00)
    {Lineárna projekcia\\[-2pt]
     \scriptsize súradnice → vnútorná reprezentácia};

\node[attnbox] (attn) at (0, -3.55)
    {Mechanizmus vlastnej pozornosti\\[-2pt]
     \scriptsize každý vrchol "sleduje" všetky ostatné};

\node[attnbox] (repeat) at (0, -5.10)
    {\(\cdots\) opakuje sa \(L\)-krát};

\node[outbox] (ctx) at (0, -6.55)
    {Kontextové reprezentácie \(z_i\)\\[-2pt]
     \scriptsize vrchol = poloha + vzťah k ostatným};

\draw[arrow] (coords) -- (embed);
\draw[arrow] (embed)  -- (attn);
\draw[arrow] (attn)   -- (repeat);
\draw[arrow] (repeat) -- (ctx);

%% Ohraničenie fázy 1
\draw[draw=orange!35, rounded corners=6pt, dashed]
    (-2.5, 0.20) rectangle (2.5, -7.20);

%% ══════════════════════════════════════════════════════════════════
%% FÁZA 2 — Skórovanie hrán
%% ══════════════════════════════════════════════════════════════════
\node[font=\bfseries\small, text=purple!70!black]
    at (7.0, 0.55) {Fáza 2 — skórovanie hrán};

%% Vstupné zdroje pre hranu {i,j}
\node[outbox] (zi) at (5.5, -0.45)
    {\(z_i\) — kontext vrcholu \(i\)};
\node[outbox] (zj) at (8.5, -0.45)
    {\(z_j\) — kontext vrcholu \(j\)};
\node[inputbox] (phi) at (7.0, -2.00)
    {\(\phi_{ij}\) — príznaky hrany};

%% Spojenie
\node[mergebox] (merge) at (7.0, -3.55)
    {Zlúčenie informácií\\[-2pt]
     \scriptsize kontext + príznaky hrany};

\node[mergebox] (head) at (7.0, -5.10)
    {Výstupná sieť\\[-2pt]
     \scriptsize (dve lineárne vrstvy + GELU)};

\node[outbox] (score) at (7.0, -6.55)
    {Skóre hrany \(s_{ij} \in \mathbb{R}\)\\[-2pt]
     \scriptsize miera príslušnosti k trase};

\draw[arrow] (zi)    |- (merge);
\draw[arrow] (zj)    |- (merge);
\draw[arrow] (phi)   -- (merge);
\draw[arrow] (merge) -- (head);
\draw[arrow] (head)  -- (score);

%% Prepojenie fázy 1 → fáza 2
\draw[arrow, draw=orange!60] (ctx.east) -- ++(1.0,0)
    |- node[midway, above, font=\scriptsize, text=orange!70]
       {pre každú hranu} (zi.west);
\draw[arrow, draw=orange!60] (ctx.east) -- ++(1.0,0)
    |- (zj.west);

%% Ohraničenie fázy 2
\draw[draw=purple!30, rounded corners=6pt, dashed]
    (4.2, 0.20) rectangle (9.8, -7.20);

%% ── Vysvetlivka dole ─────────────────────────────────────────────
\node[font=\scriptsize, text=gray!65, align=center, text width=13cm]
    at (3.5, -8.10)
    {Kľúčový rozdiel oproti MLP: transformer najprv zostaví
     \emph{globálny kontext} celej inštancie (Fáza 1) a až potom
     hodnotí každú hranu — s vedomosťou o jej vzťahu ku všetkým
     ostatným vrcholom (Fáza 2).};

\end{tikzpicture}}
\begin{tikzpicture}[
    font=\small,
    box/.style={
        draw, rounded corners=5pt, align=center,
        minimum width=3.6cm, minimum height=0.85cm,
        inner sep=5pt
    },
    inputbox/.style={box, fill=gray!10, draw=gray!50},
    attnbox/.style={box, fill=orange!10, draw=orange!50},
    mergebox/.style={box, fill=purple!8, draw=purple!40},
    outbox/.style={box, fill=green!8, draw=green!50},
    labelbox/.style={font=\scriptsize\itshape, text=gray!65, align=left,
                     text width=5.0cm},
    arrow/.style={-{Latex[length=2.5mm]}, thick},
    dasharrow/.style={-{Latex[length=2mm]}, dashed, gray!60}
]

%% ══════════════════════════════════════════════════════════════════
%% FÁZA 1 — Kódovanie vrcholov
%% ══════════════════════════════════════════════════════════════════
\node[font=\bfseries\small, text=orange!70!black]
    at (0, 0.55) {Fáza 1 — kódovanie vrcholov};

\node[inputbox] (coords) at (0, -0.45)
    {Súradnice miest \((u_i, v_i)\)\\[-2pt]
     \scriptsize pre všetky vrcholy inštancie};

\node[attnbox] (embed) at (0, -2.00)
    {Lineárna projekcia\\[-2pt]
     \scriptsize súradnice → vnútorná reprezentácia};

\node[attnbox] (attn) at (0, -3.55)
    {Mechanizmus vlastnej pozornosti\\[-2pt]
     \scriptsize každý vrchol "sleduje" všetky ostatné};

\node[attnbox] (repeat) at (0, -5.10)
    {\(\cdots\) opakuje sa \(L\)-krát};

\node[outbox] (ctx) at (0, -6.55)
    {Kontextové reprezentácie \(z_i\)\\[-2pt]
     \scriptsize vrchol = poloha + vzťah k ostatným};

\draw[arrow] (coords) -- (embed);
\draw[arrow] (embed)  -- (attn);
\draw[arrow] (attn)   -- (repeat);
\draw[arrow] (repeat) -- (ctx);

%% Ohraničenie fázy 1
\draw[draw=orange!35, rounded corners=6pt, dashed]
    (-2.5, 0.20) rectangle (2.5, -7.20);

%% ══════════════════════════════════════════════════════════════════
%% FÁZA 2 — Skórovanie hrán
%% ══════════════════════════════════════════════════════════════════
\node[font=\bfseries\small, text=purple!70!black]
    at (7.0, 0.55) {Fáza 2 — skórovanie hrán};

%% Vstupné zdroje pre hranu {i,j}
\node[outbox] (zi) at (5.5, -0.45)
    {\(z_i\) — kontext vrcholu \(i\)};
\node[outbox] (zj) at (8.5, -0.45)
    {\(z_j\) — kontext vrcholu \(j\)};
\node[inputbox] (phi) at (7.0, -2.00)
    {\(\phi_{ij}\) — príznaky hrany};

%% Spojenie
\node[mergebox] (merge) at (7.0, -3.55)
    {Zlúčenie informácií\\[-2pt]
     \scriptsize kontext + príznaky hrany};

\node[mergebox] (head) at (7.0, -5.10)
    {Výstupná sieť\\[-2pt]
     \scriptsize (dve lineárne vrstvy + GELU)};

\node[outbox] (score) at (7.0, -6.55)
    {Skóre hrany \(s_{ij} \in \mathbb{R}\)\\[-2pt]
     \scriptsize miera príslušnosti k trase};

\draw[arrow] (zi)    |- (merge);
\draw[arrow] (zj)    |- (merge);
\draw[arrow] (phi)   -- (merge);
\draw[arrow] (merge) -- (head);
\draw[arrow] (head)  -- (score);

%% Prepojenie fázy 1 → fáza 2
\draw[arrow, draw=orange!60] (ctx.east) -- ++(1.0,0)
    |- node[midway, above, font=\scriptsize, text=orange!70]
       {pre každú hranu} (zi.west);
\draw[arrow, draw=orange!60] (ctx.east) -- ++(1.0,0)
    |- (zj.west);

%% Ohraničenie fázy 2
\draw[draw=purple!30, rounded corners=6pt, dashed]
    (4.2, 0.20) rectangle (9.8, -7.20);

%% ── Vysvetlivka dole ─────────────────────────────────────────────
\node[font=\scriptsize, text=gray!65, align=center, text width=13cm]
    at (3.5, -8.10)
    {Kľúčový rozdiel oproti MLP: transformer najprv zostaví
     \emph{globálny kontext} celej inštancie (Fáza 1) a až potom
     hodnotí každú hranu — s vedomosťou o jej vzťahu ku všetkým
     ostatným vrcholom (Fáza 2).};

\end{tikzpicture}}
\begin{tikzpicture}[
    font=\small,
    box/.style={
        draw, rounded corners=5pt, align=center,
        minimum width=3.6cm, minimum height=0.85cm,
        inner sep=5pt
    },
    inputbox/.style={box, fill=gray!10, draw=gray!50},
    attnbox/.style={box, fill=orange!10, draw=orange!50},
    mergebox/.style={box, fill=purple!8, draw=purple!40},
    outbox/.style={box, fill=green!8, draw=green!50},
    labelbox/.style={font=\scriptsize\itshape, text=gray!65, align=left,
                     text width=5.0cm},
    arrow/.style={-{Latex[length=2.5mm]}, thick},
    dasharrow/.style={-{Latex[length=2mm]}, dashed, gray!60}
]

%% ══════════════════════════════════════════════════════════════════
%% FÁZA 1 — Kódovanie vrcholov
%% ══════════════════════════════════════════════════════════════════
\node[font=\bfseries\small, text=orange!70!black]
    at (0, 0.55) {Fáza 1 — kódovanie vrcholov};

\node[inputbox] (coords) at (0, -0.45)
    {Súradnice miest \((u_i, v_i)\)\\[-2pt]
     \scriptsize pre všetky vrcholy inštancie};

\node[attnbox] (embed) at (0, -2.00)
    {Lineárna projekcia\\[-2pt]
     \scriptsize súradnice → vnútorná reprezentácia};

\node[attnbox] (attn) at (0, -3.55)
    {Mechanizmus vlastnej pozornosti\\[-2pt]
     \scriptsize každý vrchol "sleduje" všetky ostatné};

\node[attnbox] (repeat) at (0, -5.10)
    {\(\cdots\) opakuje sa \(L\)-krát};

\node[outbox] (ctx) at (0, -6.55)
    {Kontextové reprezentácie \(z_i\)\\[-2pt]
     \scriptsize vrchol = poloha + vzťah k ostatným};

\draw[arrow] (coords) -- (embed);
\draw[arrow] (embed)  -- (attn);
\draw[arrow] (attn)   -- (repeat);
\draw[arrow] (repeat) -- (ctx);

%% Ohraničenie fázy 1
\draw[draw=orange!35, rounded corners=6pt, dashed]
    (-2.5, 0.20) rectangle (2.5, -7.20);

%% ══════════════════════════════════════════════════════════════════
%% FÁZA 2 — Skórovanie hrán
%% ══════════════════════════════════════════════════════════════════
\node[font=\bfseries\small, text=purple!70!black]
    at (7.0, 0.55) {Fáza 2 — skórovanie hrán};

%% Vstupné zdroje pre hranu {i,j}
\node[outbox] (zi) at (5.5, -0.45)
    {\(z_i\) — kontext vrcholu \(i\)};
\node[outbox] (zj) at (8.5, -0.45)
    {\(z_j\) — kontext vrcholu \(j\)};
\node[inputbox] (phi) at (7.0, -2.00)
    {\(\phi_{ij}\) — príznaky hrany};

%% Spojenie
\node[mergebox] (merge) at (7.0, -3.55)
    {Zlúčenie informácií\\[-2pt]
     \scriptsize kontext + príznaky hrany};

\node[mergebox] (head) at (7.0, -5.10)
    {Výstupná sieť\\[-2pt]
     \scriptsize (dve lineárne vrstvy + GELU)};

\node[outbox] (score) at (7.0, -6.55)
    {Skóre hrany \(s_{ij} \in \mathbb{R}\)\\[-2pt]
     \scriptsize miera príslušnosti k trase};

\draw[arrow] (zi)    |- (merge);
\draw[arrow] (zj)    |- (merge);
\draw[arrow] (phi)   -- (merge);
\draw[arrow] (merge) -- (head);
\draw[arrow] (head)  -- (score);

%% Prepojenie fázy 1 → fáza 2
\draw[arrow, draw=orange!60] (ctx.east) -- ++(1.0,0)
    |- node[midway, above, font=\scriptsize, text=orange!70]
       {pre každú hranu} (zi.west);
\draw[arrow, draw=orange!60] (ctx.east) -- ++(1.0,0)
    |- (zj.west);

%% Ohraničenie fázy 2
\draw[draw=purple!30, rounded corners=6pt, dashed]
    (4.2, 0.20) rectangle (9.8, -7.20);

%% ── Vysvetlivka dole ─────────────────────────────────────────────
\node[font=\scriptsize, text=gray!65, align=center, text width=13cm]
    at (3.5, -8.10)
    {Kľúčový rozdiel oproti MLP: transformer najprv zostaví
     \emph{globálny kontext} celej inštancie (Fáza 1) a až potom
     hodnotí každú hranu — s vedomosťou o jej vzťahu ku všetkým
     ostatným vrcholom (Fáza 2).};

\end{tikzpicture}}
\begin{tikzpicture}[
    font=\small,
    box/.style={
        draw, rounded corners=5pt, align=center,
        minimum width=3.6cm, minimum height=0.85cm,
        inner sep=5pt
    },
    inputbox/.style={box, fill=gray!10, draw=gray!50},
    attnbox/.style={box, fill=orange!10, draw=orange!50},
    mergebox/.style={box, fill=purple!8, draw=purple!40},
    outbox/.style={box, fill=green!8, draw=green!50},
    labelbox/.style={font=\scriptsize\itshape, text=gray!65, align=left,
                     text width=5.0cm},
    arrow/.style={-{Latex[length=2.5mm]}, thick},
    dasharrow/.style={-{Latex[length=2mm]}, dashed, gray!60}
]

%% ══════════════════════════════════════════════════════════════════
%% FÁZA 1 — Kódovanie vrcholov
%% ══════════════════════════════════════════════════════════════════
\node[font=\bfseries\small, text=orange!70!black]
    at (0, 0.55) {Fáza 1 — kódovanie vrcholov};

\node[inputbox] (coords) at (0, -0.45)
    {Súradnice miest \((u_i, v_i)\)\\[-2pt]
     \scriptsize pre všetky vrcholy inštancie};

\node[attnbox] (embed) at (0, -2.00)
    {Lineárna projekcia\\[-2pt]
     \scriptsize súradnice → vnútorná reprezentácia};

\node[attnbox] (attn) at (0, -3.55)
    {Mechanizmus vlastnej pozornosti\\[-2pt]
     \scriptsize každý vrchol "sleduje" všetky ostatné};

\node[attnbox] (repeat) at (0, -5.10)
    {\(\cdots\) opakuje sa \(L\)-krát};

\node[outbox] (ctx) at (0, -6.55)
    {Kontextové reprezentácie \(z_i\)\\[-2pt]
     \scriptsize vrchol = poloha + vzťah k ostatným};

\draw[arrow] (coords) -- (embed);
\draw[arrow] (embed)  -- (attn);
\draw[arrow] (attn)   -- (repeat);
\draw[arrow] (repeat) -- (ctx);

%% Ohraničenie fázy 1
\draw[draw=orange!35, rounded corners=6pt, dashed]
    (-2.5, 0.20) rectangle (2.5, -7.20);

%% ══════════════════════════════════════════════════════════════════
%% FÁZA 2 — Skórovanie hrán
%% ══════════════════════════════════════════════════════════════════
\node[font=\bfseries\small, text=purple!70!black]
    at (7.0, 0.55) {Fáza 2 — skórovanie hrán};

%% Vstupné zdroje pre hranu {i,j}
\node[outbox] (zi) at (5.5, -0.45)
    {\(z_i\) — kontext vrcholu \(i\)};
\node[outbox] (zj) at (8.5, -0.45)
    {\(z_j\) — kontext vrcholu \(j\)};
\node[inputbox] (phi) at (7.0, -2.00)
    {\(\phi_{ij}\) — príznaky hrany};

%% Spojenie
\node[mergebox] (merge) at (7.0, -3.55)
    {Zlúčenie informácií\\[-2pt]
     \scriptsize kontext + príznaky hrany};

\node[mergebox] (head) at (7.0, -5.10)
    {Výstupná sieť\\[-2pt]
     \scriptsize (dve lineárne vrstvy + GELU)};

\node[outbox] (score) at (7.0, -6.55)
    {Skóre hrany \(s_{ij} \in \mathbb{R}\)\\[-2pt]
     \scriptsize miera príslušnosti k trase};

\draw[arrow] (zi)    |- (merge);
\draw[arrow] (zj)    |- (merge);
\draw[arrow] (phi)   -- (merge);
\draw[arrow] (merge) -- (head);
\draw[arrow] (head)  -- (score);

%% Prepojenie fázy 1 → fáza 2
\draw[arrow, draw=orange!60] (ctx.east) -- ++(1.0,0)
    |- node[midway, above, font=\scriptsize, text=orange!70]
       {pre každú hranu} (zi.west);
\draw[arrow, draw=orange!60] (ctx.east) -- ++(1.0,0)
    |- (zj.west);

%% Ohraničenie fázy 2
\draw[draw=purple!30, rounded corners=6pt, dashed]
    (4.2, 0.20) rectangle (9.8, -7.20);

%% ── Vysvetlivka dole ─────────────────────────────────────────────
\node[font=\scriptsize, text=gray!65, align=center, text width=13cm]
    at (3.5, -8.10)
    {Kľúčový rozdiel oproti MLP: transformer najprv zostaví
     \emph{globálny kontext} celej inštancie (Fáza 1) a až potom
     hodnotí každú hranu — s vedomosťou o jej vzťahu ku všetkým
     ostatným vrcholom (Fáza 2).};

\end{tikzpicture}